%! TeX program = lualatex
% Nineteen pages: the whole map, then nine layers of it -- the two-sided
% limit branch, the first decision, the 0/0 case, the nonzero/0 case, the
% evaluation techniques, the continuity branch, the continuity test, why
% substitution works, the discontinuous branch -- each set twice, once on
% page 1's canvas and once trimmed to its own content.
% The quarter-turn print variant of page 1 lives in
% mindmap-limits-bigger.tex.  `multi' is what
% makes each mmpage a page of its own -- the [tikz] option would do that for a
% bare tikzpicture, but it leaves the picture ungrouped and \resizebox below
% then fails with a missing \endgroup.
\documentclass[multi=mmpage]{standalone}
\input{../typography.tex.preamble}

% DO NOT use the tikz preamble: its `basic curve' styles pull in pgfplots, and
% enabling `trig format=rad' there makes this picture die with a dimension too
% large error.
% \input{../tikz.tex.preamble}
\usepackage{graphicx}
\usepackage{tikz}

\usetikzlibrary{mindmap, backgrounds, calc}

% The "substitution just works" case.  WesternSpring on its own is a neon
% lime; muting it through grey keeps it in the palette without shouting.
\colorlet{substcase}{WesternSpring!55!gray!45}
% Same hue at full strength, for the arrow joining two substcase nodes.
\colorlet{substline}{WesternSpring!55!gray}
% Both one-sided limits exist.  WesternSky neat is a neon cyan, so mute it
% through grey the same way.
\colorlet{bothexist}{WesternSky!60!gray!45}
% Connector labels.  Neutral and italic rather than coloured: the node fills
% already carry the colour coding, and `attn' gold went muddy beside them.
\colorlet{decision}{black!60}

\newenvironment{mmpage}{}{}

% Every page after the first is a layer of page 1, and each layer is set
% twice.  The first of the pair uses this canvas -- page 1's own bounding
% box, measured from the root and rounded outward to whole points -- so the
% page comes out page 1's size with every node where page 1 puts it, and
% the set flips through as overlays.  The second of the pair leaves the
% canvas off, so standalone trims the page to the layer's own content; its
% \resizebox width is page 1's scale factor (7in over its natural 815.6pt)
% applied to that trimmed width, which makes it the same picture with the
% empty space cut away rather than a blown-up one.
\newcommand{\pagecanvas}{%
  \path [use as bounding box] (-409pt,-335pt) rectangle (407pt,351pt);}

% Second line of a concept label.  Formulas stay in \displaystyle, as
% \everymath sets course-wide, so \lim subscripts sit underneath: the circles
% below are sized to allow for that.
\newcommand{\sub}[1]{\\[0.3ex]{\scriptsize\(#1\)}}
\newcommand{\note}[1]{\\[0.1ex]{\tiny #1}}

% The mindmap library gives a level-2 concept 1.5cm
% of text width and a level-3 one just 1cm, which hyphenates single words like
% remov-able and Inter-mediate.  Widen them, and keep the manual line breaks
% short enough to fit what is set here.
\tikzset{
  mindmap base/.style={
    mindmap,
    grow cyclic,
    concept color=gray!20,
    every child node/.style={concept, minimum size=1.4cm, text width=1.7cm,
                             align=center, font=\scriptsize},
  },
  % Connector labels.  `mindmap' hands every node a large minimum size; that is
  % harmless when a label is centred on a point, but an edge-anchored one ends
  % up centimetres off its bar unless the box is shrunk back to its text.
  barlabel/.style={font=\footnotesize\itshape, text=decision, align=center,
                   inner sep=2pt, minimum size=0pt,
                   minimum width=0pt, minimum height=0pt},
}

% The level distances and sibling angles page 1 draws with.  Every layer
% lifted from it is drawn with the same ones, so a node cannot land
% anywhere but where page 1 puts it.
\tikzset{
  map levels/.style={
    mindmap base,
    level 1/.append style={level distance=5.5cm, sibling angle=120},
    level 2/.append style={level distance=4cm, sibling angle=60},
    level 3/.append style={level distance=3.6cm, sibling angle=44},
  },
  % A layer with `dirsub' for its own root: what page 1 sets for levels 2
  % and 3 is set here for levels 1 and 2.  Split in two so a layer that
  % draws this subtree beside another one can put just the levels in a
  % scope of their own.
  technique levels/.style={
    level 1/.append style={level distance=4cm, sibling angle=60},
    level 2/.append style={level distance=4.2cm, sibling angle=40},
  },
  extract levels/.style={mindmap base, technique levels},
  % Likewise with `conttest' for the root.  The 2.1cm is what page 1's
  % continuity branch sets for everything in it, and it belongs in a scope
  % around the tree rather than on the whole picture: the connector labels
  % break their own lines and must not be given a text width.
  continuity levels/.style={
    every node/.append style={text width=2.1cm},
    level 1/.append style={level distance=5cm, sibling angle=50},
    level 2/.append style={level distance=3.6cm, sibling angle=52},
  },
}

% Two things about the angles on every page.
%
% (1) A `clockwise from' given to the root's children leaks down to every
%     deeper level, so without a correction every node's children fan in the
%     same absolute direction instead of pointing outward.  Each node
%     therefore restates its own fan, in the path right after it -- in a
%     child's option list the key would instead re-aim that child itself.
%     The value is (parent angle) - (children - 1) * (sibling angle) / 2.
%
% (2) Level distances and sibling angles differ per subtree, so they are set
%     in each child's option list, where they propagate down that subtree.

\begin{document}
% ===================================================================
%  Page 1 -- limits at a number
% ===================================================================
\begin{mmpage}
% The guided notes want a 7in-wide graphic (mindmap-integration.pdf is 7.3in).
% The picture is drawn wider than that and scaled down to fit, so keep labels
% short: at this size every extra line of text costs legibility.
\resizebox{7in}{!}{%
\begin{tikzpicture}[map levels]
  \pagecanvas
  % Branches at 135, 45, -45, -135.  Continuity sits next to direct
  % substitution so the cross-connection between them stays short.
  % The number line is a nested tikzpicture inside the root's label: no tick
  % labels, one dot at a, and the axis tip named x.
  \node [concept, minimum size=3.8cm, text width=3.2cm, font=\small] (root)
    {limit at a constant \(a\)\\[0.8ex]
     % The nested picture inherits the mindmap's node sizing, which flings
     % these two labels centimetres away; reset it here.
     \begin{tikzpicture}[baseline,
         every node/.style={inner sep=1.5pt, minimum size=0pt, minimum width=0pt,
                            minimum height=0pt, text width=, align=center}]
       \draw [->, thin] (0,0) -- (2.5,0) node [right] {\scriptsize\(x\)};
       \fill (1.25,0) circle (1.4pt);
       \node [below] at (1.25,0) {\scriptsize\(a\)};
     \end{tikzpicture}}
    [clockwise from=150]
    child [concept color=gray!20,
           every node/.append style={text width=2.1cm},
           level 2/.append style={level distance=5cm, sibling angle=50},
           level 3/.append style={level distance=3.6cm, sibling angle=52}] { 
      % Wide enough to keep the test on one line -- a node option beats the
      % 2.1cm this subtree sets for everything below it.
      node [minimum size=3.2cm, text width=3.4cm, font=\footnotesize] (conttest)
        {continuity at \(a\)\sub{\lim_{x \to a} f(x) \overset{?}{=} f(a)}}
      [counterclockwise from=75]
      child [concept color=substcase] { 
        node (yes) {continuous}
        [counterclockwise from=65]
        child [concept color=blue!20] { node (ivt) {Intermediate\\Value Theorem} }
        child { node (families) [text width=2.6cm] {can use direct\\substitution\note{to evaluate limits of\\polynomials, rational,\\trigs, \(e^{x}\), \(\ln(x)\)\\in their domain}} }
      }
      child { node (leftcont) [minimum size=2.2cm, text width=2.4cm]
        {left-continuous\note{\(\lim_{x \to a^{-}} f(x) \overset{?}{=} f(a)\)}} }
      child { node (rightcont) [minimum size=2.2cm, text width=2.4cm]
        {right-continuous\note{\(\lim_{x \to a^{+}} f(x) \overset{?}{=} f(a)\)}} }
      child [concept color=gray!20,
             level 3/.append style={level distance=5.2cm, sibling angle=40}] { 
        node (no) {discontinuous\note{calculate both\\one-sided limits}}
        [counterclockwise from=187]
        % Cyan for the two cases where both one-sided limits exist -- not gold,
        % which is the connector-label colour.  The infinite case keeps the
        % pink it shares with nonzero/0.
        child [concept color=bothexist] { node (removable) {removable} }
        child [concept color=bothexist] { node (jump) {jump} }
        child [concept color=pink!50] { node (infinite) [text width=2.6cm]
          {infinite and \(f\) has\\a vertical\\asymptote at \(a\)} }
      }
    }
    % Pushed further out than its siblings: the bar into it carries a
    % two-line label, and the techniques hanging off it need the room.
    child [concept color=gray!20, level distance=5.5cm,
           level 3/.append style={level distance=4.2cm, sibling angle=40}] { 
      node [minimum size=2.4cm, text width=2.2cm, font=\footnotesize] (dirsub)
        {try direct\\substitution\note{is the denominator \(0\)?}}
      [counterclockwise from=-30]
      child [concept color=supp!20] { 
        node (zerozero) {\(\frac{0}{0}\)}
        % Four children now, and the fan is swung well below this node's own
        % direction so the last of them clears the nonzero/0 subtree above.
        [counterclockwise from=-139]
        child { node (squeeze) {Squeeze\\Theorem} }
        child { node (cov) {change of\\variable} }
        child { node (rationalize) {rationalize\\and cancel} }
        child { node (factor) {factor and\\cancel} }
        % Its own colour, reserved for differentiation elsewhere in the notes.
        child [concept color=main!25] { node (deriv) [text width=2.2cm]
          {recognize as\\a derivative} }
      }
      % Pink to pair it with the infinite discontinuity the arrow points at.
      child [concept color=pink!50,
             level 3/.append style={level distance=5.2cm}] { 
        node (constzero) {\(\frac{\text{nonzero}}{0}\)}
        [counterclockwise from=10]
        child { node (differ) {limit does\\not exist} }
        child { node (agree) {\(+\infty\) or \(-\infty\)} }
      }
      child [concept color=substcase] { node (algno) [text width=2.1cm]
          {\(\frac{\text{any number}}{\text{nonzero}}\)} }
    }
    child [concept color=gray!20,
           level 2/.append style={level distance=5cm}] { 
      node (twosided) [minimum size=2.2cm, text width=2.3cm]
        {two-sided limit\sub{\lim_{x \to a} f(x)}}
      [counterclockwise from=-150]
      child { node (leftlim) [minimum size=2.2cm, text width=2.3cm]
        {left limit at \(a\)\sub{\lim_{x \to a^{-}} f(x)}} }
      child { node (rightlim) [minimum size=2.2cm, text width=2.3cm]
        {right limit at \(a\)\sub{\lim_{x \to a^{+}} f(x)}} }
      child [concept color=supp!20] { node (known) [text width=2cm]
        {well-known\\limits} }
    }
    ;

  \node [concept color=white, annotation, text width=, right,
         font=\scriptsize, align=center] at (known.east)
    % `aligned' rather than `align*': a display environment inside a TikZ node
    % dies with a missing \endgroup.  \everymath is \displaystyle, so this
    % sets exactly as align* would.
    {\(\begin{aligned}
       \lim_{x \to 0} \frac{\sin(x)}{x} &= 1\\
       \lim_{x \to 0} \frac{\cos(x) - 1}{x} &= 0\\
       \lim_{x \to 0} \frac{e^{x} - 1}{x} &= 1\\
       \lim_{x \to \infty} \left(1 + \frac{1}{x}\right)^{x} &= e
     \end{aligned}\)};

  \node [concept color=white, annotation, right, text width=4.2cm,
         font=\scriptsize, align=left] at (ivt.east)
    {IVT applies to continuous function over \([a,b]\) and \(f(a) \ne f(b)\)};

  % Labels riding the bars into a branch, rather than sitting in it as one
  % more node.  All of them are decisions or instructions, so they share the
  % `attn' colour and read as a layer distinct from the concept names.
  % In the foreground, not on the arrow's own path: the arrow is drawn on the
  % background layer, so a label riding it disappears under the nodes it passes.
  \node [barlabel, align=left, anchor=west]
    at ([shift={(0.35cm,0)}]$(cov)!0.5!(known)$) {in some cases,\\we use \dots};
  \node [barlabel] at ($(no)!0.5!(removable)$)
    {both exist and agree,\\but \(\ne f(a)\)};
  \node [barlabel] at ($(no)!0.5!(jump)$) {both exist\\but disagree};
  \node [barlabel] at ($(no)!0.5!(infinite)$)
    {one of them\\is \(\pm\infty\)};
  \node [barlabel] at ($(constzero)!0.5!(agree)$)
    {one-sided\\limits agree};
  \node [barlabel] at ($(constzero)!0.5!(differ)$)
    {one-sided\\limits disagree};
  \node [barlabel] at ($(conttest)!0.5!(yes)$) {yes};
  \node [barlabel] at ($(conttest)!0.5!(no)$) {no};
  \node [barlabel] at ($(root)!0.5!(dirsub)$) {calculations uses\\limit laws};
  % Graphs lifted from plot-limit-intro.pdf and parked beside the node each
  % one illustrates.  \includegraphics reads the built PDF, so the path is
  % relative to this directory -- and plot-limit-intro.pdf has to exist before
  % this file is compiled.
  % Anchored off the node each one illustrates, so they follow if the layout
  % moves.  `annotation' forces a 3.5cm ragged text box and these images are
  % 3.81cm wide, so reset the text width to let each node size to its image --
  % otherwise every one of them overfulls and hangs past its own node box.
  \node [concept color=white, annotation, text width=, right]
    at (twosided.east) {\includegraphics[width=0.75in, page=3]{plot-limit-intro}};
  \node [concept color=white, annotation, text width=, below]
    at (leftlim.south) {\includegraphics[width=0.75in, page=1]{plot-limit-intro}};
  \node [concept color=white, annotation, text width=, right]
    at (rightlim.east) {\includegraphics[width=0.75in, page=2]{plot-limit-intro}};

  % The one-sided continuity pictures, smaller because there are four of them
  % in a crowded corner.  The last is stacked above its sibling rather than
  % anchored straight off the node: everything below rightcont is taken.
  \node [concept color=white, annotation, text width=, above]
    at (leftcont.north) {\includegraphics[width=0.75in, page=13]{plot-limit-intro}};
  \node [concept color=white, annotation, text width=, above left]
    at (leftcont.west) {\includegraphics[width=0.75in, page=14]{plot-limit-intro}};
  \node [concept color=white, annotation, text width=, left]
    at (rightcont.west) {\includegraphics[width=0.75in, page=15]{plot-limit-intro}};
  \node [concept color=white, annotation, text width=, left]
    at ([shift={(0,2.1cm)}]rightcont.west) {\includegraphics[width=0.75in, page=16]{plot-limit-intro}};



  % Nudged clear of the connection it labels -- the label sits on the main
  % layer and the bar below it is saturated, so text straight on top of it
  % would not read.
  \node [barlabel, anchor=south]
    at ($(conttest)!0.5!(dirsub)$) {testing for continuity\\uses these techniques};

  % The two decision trees reach the same place by different routes.
  \begin{pgfonlayer}{background}
    \draw [concept connection, substline, ->, >=stealth, shorten >=4pt]
      (algno) edge (yes);
    % Directed: continuity is what borrows the techniques, not the other way
    % round.  Shortened so the head clears the concept's blob outline.  Drawn
    % in the root's own grey, so it reads as part of the map's spine rather
    % than as one of the coloured cross-links.
    \draw [concept connection, gray!20, ->, >=stealth, shorten >=4pt]
      (conttest) edge (dirsub);
    \draw [concept connection, supp!60, ->, >=stealth, shorten >=4pt]
      (cov) edge (known);
    % A circular arc, bowed away from the root so it threads the gap between
    % the root and its branches.  No heads, to match the other plain link.
    \draw [concept connection, pink]
      (constzero.south) to [bend left=20] (infinite);
  \end{pgfonlayer}
\end{tikzpicture} 
}
\end{mmpage}

% ===================================================================
%  Pages 2 and 3 -- the two-sided limit layer: the root, the branch that defines
%  the limit, and the graphs and well-known limits hanging off it.
%
%  Only the root's third branch is drawn, so the fan starts at that
%  branch's own angle instead of at 150: with a single child, `clockwise
%  from=150' would put it where the first branch goes.  Everything else is
%  page 1's code unchanged.
%
%  Page 2 is the layer on page 1's canvas, page 3 the same drawing
%  trimmed to its own content.
% ===================================================================
\newcommand{\twosidedlayer}{%
  \node [concept, minimum size=3.8cm, text width=3.2cm, font=\small] (root)
    {limit at a constant \(a\)\\[0.8ex]
     % The nested picture inherits the mindmap's node sizing, which flings
     % these two labels centimetres away; reset it here.
     \begin{tikzpicture}[baseline,
         every node/.style={inner sep=1.5pt, minimum size=0pt, minimum width=0pt,
                            minimum height=0pt, text width=, align=center}]
       \draw [->, thin] (0,0) -- (2.5,0) node [right] {\scriptsize\(x\)};
       \fill (1.25,0) circle (1.4pt);
       \node [below] at (1.25,0) {\scriptsize\(a\)};
     \end{tikzpicture}}
    [clockwise from=-90]
    child [concept color=gray!20,
           level 2/.append style={level distance=5cm}] { 
      node (twosided) [minimum size=2.2cm, text width=2.3cm]
        {two-sided limit\sub{\lim_{x \to a} f(x)}}
      [counterclockwise from=-150]
      child { node (leftlim) [minimum size=2.2cm, text width=2.3cm]
        {left limit at \(a\)\sub{\lim_{x \to a^{-}} f(x)}} }
      child { node (rightlim) [minimum size=2.2cm, text width=2.3cm]
        {right limit at \(a\)\sub{\lim_{x \to a^{+}} f(x)}} }
      child [concept color=supp!20] { node (known) [text width=2cm]
        {well-known\\limits} }
    }
    ;

  \node [concept color=white, annotation, text width=, right,
         font=\scriptsize, align=center] at (known.east)
    % `aligned' rather than `align*': a display environment inside a TikZ node
    % dies with a missing \endgroup.  \everymath is \displaystyle, so this
    % sets exactly as align* would.
    {\(\begin{aligned}
       \lim_{x \to 0} \frac{\sin(x)}{x} &= 1\\
       \lim_{x \to 0} \frac{\cos(x) - 1}{x} &= 0\\
       \lim_{x \to 0} \frac{e^{x} - 1}{x} &= 1\\
       \lim_{x \to \infty} \left(1 + \frac{1}{x}\right)^{x} &= e
     \end{aligned}\)};

  % Graphs lifted from plot-limit-intro.pdf and parked beside the node each
  % one illustrates.  \includegraphics reads the built PDF, so the path is
  % relative to this directory -- and plot-limit-intro.pdf has to exist before
  % this file is compiled.
  % Anchored off the node each one illustrates, so they follow if the layout
  % moves.  `annotation' forces a 3.5cm ragged text box and these images are
  % 3.81cm wide, so reset the text width to let each node size to its image --
  % otherwise every one of them overfulls and hangs past its own node box.
  \node [concept color=white, annotation, text width=, right]
    at (twosided.east) {\includegraphics[width=0.75in, page=3]{plot-limit-intro}};
  \node [concept color=white, annotation, text width=, below]
    at (leftlim.south) {\includegraphics[width=0.75in, page=1]{plot-limit-intro}};
  \node [concept color=white, annotation, text width=, right]
    at (rightlim.east) {\includegraphics[width=0.75in, page=2]{plot-limit-intro}};
}

\begin{mmpage}
\resizebox{7in}{!}{%
\begin{tikzpicture}[map levels]
  \pagecanvas
  \twosidedlayer
\end{tikzpicture} 
}
\end{mmpage}

% The same layer, trimmed to its own content.
\begin{mmpage}
\resizebox{3.292in}{!}{%
\begin{tikzpicture}[map levels]
  \twosidedlayer
\end{tikzpicture} 
}
\end{mmpage}

% ===================================================================
%  Pages 4 and 5 -- the first decision: the root, `try direct substitution', and
%  the three cases the denominator can fall into.  What to do about each
%  is page 4.
%
%  Same one-branch trick as page 2, aimed at the second branch (30).  The
%  three cases keep the colours and the fan page 1 gives them; only their
%  own children are left out.
%
%  Page 4 is the layer on page 1's canvas, page 5 the same drawing
%  trimmed to its own content.
% ===================================================================
\newcommand{\firstdecisionlayer}{%
  \node [concept, minimum size=3.8cm, text width=3.2cm, font=\small] (root)
    {limit at a constant \(a\)\\[0.8ex]
     % The nested picture inherits the mindmap's node sizing, which flings
     % these two labels centimetres away; reset it here.
     \begin{tikzpicture}[baseline,
         every node/.style={inner sep=1.5pt, minimum size=0pt, minimum width=0pt,
                            minimum height=0pt, text width=, align=center}]
       \draw [->, thin] (0,0) -- (2.5,0) node [right] {\scriptsize\(x\)};
       \fill (1.25,0) circle (1.4pt);
       \node [below] at (1.25,0) {\scriptsize\(a\)};
     \end{tikzpicture}}
    [clockwise from=30]
    child [concept color=gray!20, level distance=5.5cm,
           level 3/.append style={level distance=4.2cm, sibling angle=40}] { 
      node [minimum size=2.4cm, text width=2.2cm, font=\footnotesize] (dirsub)
        {try direct\\substitution\note{is the denominator \(0\)?}}
      [counterclockwise from=-30]
      child [concept color=supp!20] { node (zerozero) {\(\frac{0}{0}\)} }
      child [concept color=pink!50] { node (constzero)
        {\(\frac{\text{nonzero}}{0}\)} }
      child [concept color=substcase] { node (algno) [text width=2.1cm]
          {\(\frac{\text{any number}}{\text{nonzero}}\)} }
    }
    ;

  \node [barlabel] at ($(root)!0.5!(dirsub)$) {calculations uses\\limit laws};
}

\begin{mmpage}
\resizebox{7in}{!}{%
\begin{tikzpicture}[map levels]
  \pagecanvas
  \firstdecisionlayer
\end{tikzpicture} 
}
\end{mmpage}

% The same layer, trimmed to its own content.
\begin{mmpage}
\resizebox{2.707in}{!}{%
\begin{tikzpicture}[map levels]
  \firstdecisionlayer
\end{tikzpicture} 
}
\end{mmpage}


% ===================================================================
%  Pages 6 and 7 -- the 0/0 layer: `try direct substitution' and the one
%  case that needs work, with the five ways of doing it and the link from
%  the change of variable to the well-known limits.  The other two cases,
%  and the list of those limits, arrive on pages 8 and 9.
%
%  Drawn like the techniques layer, with `dirsub' as its own root at the
%  point page 1 puts it, and the same fan for the 0/0 subtree, so both
%  pages register with page 1.
%
%  Page 6 is the layer on page 1's canvas, page 7 the same drawing
%  trimmed to its own content.
% ===================================================================
\newcommand{\zerozerolayer}{%
  \node [concept, minimum size=2.4cm, text width=2.2cm, align=center,
         font=\footnotesize] (dirsub) at (30:5.5cm)
    {try direct\\substitution\note{is the denominator \(0\)?}}
    [counterclockwise from=-30]
    child [concept color=supp!20] { 
      node (zerozero) {\(\frac{0}{0}\)}
      % Four children now, and the fan is swung well below this node's own
      % direction so the last of them clears the nonzero/0 subtree above.
      [counterclockwise from=-139]
      child { node (squeeze) {Squeeze\\Theorem} }
      child { node (cov) {change of\\variable} }
      child { node (rationalize) {rationalize\\and cancel} }
      child { node (factor) {factor and\\cancel} }
      % Its own colour, reserved for differentiation elsewhere in the notes.
      child [concept color=main!25] { node (deriv) [text width=2.2cm]
        {recognize as\\a derivative} }
    }
    ;

  % Off the tree: on page 1 this hangs from the two-sided branch, which is
  % not part of this layer, so it is placed at the point that branch would
  % put it and drawn with the styling `every child node' would have given it.
  % Its list of well-known limits is left for the layer that follows.
  \node [concept color=supp!20, concept, minimum size=1.4cm, text width=2cm,
         align=center, font=\scriptsize] (known)
    at ($(-90:5.5cm)+(-30:5cm)$) {well-known\\limits};

  % In the foreground, not on the arrow's own path: the arrow is drawn on
  % the background layer, so a label riding it disappears under the nodes it
  % passes.
  \node [barlabel, align=left, anchor=west]
    at ([shift={(0.35cm,0)}]$(cov)!0.5!(known)$) {in some cases,\\we use \dots};

  \begin{pgfonlayer}{background}
    \draw [concept connection, supp!60, ->, >=stealth, shorten >=4pt]
      (cov) edge (known);
  \end{pgfonlayer}
}

\begin{mmpage}
\resizebox{7in}{!}{%
\begin{tikzpicture}[extract levels]
  \pagecanvas
  \zerozerolayer
\end{tikzpicture} 
}
\end{mmpage}

% The same layer, trimmed to its own content.
\begin{mmpage}
\resizebox{2.464in}{!}{%
\begin{tikzpicture}[extract levels]
  \zerozerolayer
\end{tikzpicture} 
}
\end{mmpage}

% ===================================================================
%  Pages 8 and 9 -- the nonzero/0 layer: `try direct substitution' with all
%  three cases, and the one that ends in an infinity worked out.  The 0/0
%  case keeps its place in the fan but not its own subtree, which is pages
%  6 and 7.
%
%  Drawn like the other techniques layers, with `dirsub' as its own root at
%  the point page 1 puts it, so both pages register with page 1.
%
%  Page 8 is the layer on page 1's canvas, page 9 the same drawing trimmed
%  to its own content.
% ===================================================================
\newcommand{\nonzerolayer}{%
  \node [concept, minimum size=2.4cm, text width=2.2cm, align=center,
         font=\footnotesize] (dirsub) at (30:5.5cm)
    {try direct\\substitution\note{is the denominator \(0\)?}}
    [counterclockwise from=-30]
    child [concept color=supp!20] { node (zerozero) {\(\frac{0}{0}\)} }
    % Pink on page 1 to pair it with the infinite discontinuity its arc
    % points at.  That arc is not in this layer, but the colour stays.
    child [concept color=pink!50,
           level 2/.append style={level distance=5.2cm}] { 
      node (constzero) {\(\frac{\text{nonzero}}{0}\)}
      [counterclockwise from=10]
      child { node (differ) {limit does\\not exist} }
      child { node (agree) {\(+\infty\) or \(-\infty\)} }
    }
    child [concept color=substcase] { node (algno) [text width=2.1cm]
        {\(\frac{\text{any number}}{\text{nonzero}}\)} }
    ;

  \node [barlabel] at ($(constzero)!0.5!(agree)$)
    {one-sided\\limits agree};
  \node [barlabel] at ($(constzero)!0.5!(differ)$)
    {one-sided\\limits disagree};
}

\begin{mmpage}
\resizebox{7in}{!}{%
\begin{tikzpicture}[extract levels]
  \pagecanvas
  \nonzerolayer
\end{tikzpicture} 
}
\end{mmpage}

% The same layer, trimmed to its own content.
\begin{mmpage}
\resizebox{2.633in}{!}{%
\begin{tikzpicture}[extract levels]
  \nonzerolayer
\end{tikzpicture} 
}
\end{mmpage}

% ===================================================================
%  Pages 10 and 11 -- the evaluation techniques on their own: the `try direct
%  substitution' subtree of page 1, plus the well-known limits it feeds
%  into, lifted out of the full map.
%
%  Every node keeps the position page 1 gives it, measured from that
%  page's root, so this layer registers with page 1 like the others do.
%  That is why the subtree is not simply re-rooted: `dirsub' is placed
%  by hand at the root's second branch (clockwise from 150, so 30 degrees
%  out, at the 5.5cm level-1 distance), and `known' at the end of the
%  two-sided branch (-90 for 5.5cm, then -30 for the 5cm that branch
%  sets for its own children).  The root itself is not drawn here.
%
%  Levels shift up by one once `dirsub' is the root, so what page 1 sets
%  for levels 2 and 3 is set here for levels 1 and 2.  Node sizes do not
%  shift with them: `every child node' in `mindmap base' pins those for
%  every level alike.
%
%  Page 10 is the layer on page 1's canvas, page 11 the same drawing
%  trimmed to its own content.
% ===================================================================
\newcommand{\techniqueslayer}{%
  \node [concept, minimum size=2.4cm, text width=2.2cm, align=center,
         font=\footnotesize] (dirsub) at (30:5.5cm)
    {try direct\\substitution\note{is the denominator \(0\)?}}
    [counterclockwise from=-30]
    child [concept color=supp!20] { 
      node (zerozero) {\(\frac{0}{0}\)}
      % Four children now, and the fan is swung well below this node's own
      % direction so the last of them clears the nonzero/0 subtree above.
      [counterclockwise from=-139]
      child { node (squeeze) {Squeeze\\Theorem} }
      child { node (cov) {change of\\variable} }
      child { node (rationalize) {rationalize\\and cancel} }
      child { node (factor) {factor and\\cancel} }
      % Its own colour, reserved for differentiation elsewhere in the notes.
      child [concept color=main!25] { node (deriv) [text width=2.2cm]
        {recognize as\\a derivative} }
    }
    % Pink on page 1 to pair it with the infinite discontinuity its arc
    % points at.  That arc is not in this extract, but the colour stays:
    % the two drawings are meant to be read against each other.
    child [concept color=pink!50,
           level 2/.append style={level distance=5.2cm}] { 
      node (constzero) {\(\frac{\text{nonzero}}{0}\)}
      [counterclockwise from=10]
      child { node (differ) {limit does\\not exist} }
      child { node (agree) {\(+\infty\) or \(-\infty\)} }
    }
    child [concept color=substcase] { node (algno) [text width=2.1cm]
        {\(\frac{\text{any number}}{\text{nonzero}}\)} }
    ;

  % Off the tree: on page 1 this hangs from the two-sided branch, which is
  % not part of the extract, so it is placed at the point that branch would
  % put it and drawn with the styling `every child node' would have given it.
  \node [concept color=supp!20, concept, minimum size=1.4cm, text width=2cm,
         align=center, font=\scriptsize] (known)
    at ($(-90:5.5cm)+(-30:5cm)$) {well-known\\limits};

  \node [concept color=white, annotation, text width=, right,
         font=\scriptsize, align=center] at (known.east)
    % `aligned' rather than `align*': a display environment inside a TikZ node
    % dies with a missing \endgroup.  \everymath is \displaystyle, so this
    % sets exactly as align* would.
    {\(\begin{aligned}
       \lim_{x \to 0} \frac{\sin(x)}{x} &= 1\\
       \lim_{x \to 0} \frac{\cos(x) - 1}{x} &= 0\\
       \lim_{x \to 0} \frac{e^{x} - 1}{x} &= 1\\
       \lim_{x \to \infty} \left(1 + \frac{1}{x}\right)^{x} &= e
     \end{aligned}\)};

  % Labels riding the bars into a branch, rather than sitting in it as one
  % more node.  In the foreground, not on the arrow's own path: the arrow is
  % drawn on the background layer, so a label riding it disappears under the
  % nodes it passes.
  \node [barlabel, align=left, anchor=west]
    at ([shift={(0.35cm,0)}]$(cov)!0.5!(known)$) {in some cases,\\we use \dots};
  \node [barlabel] at ($(constzero)!0.5!(agree)$)
    {one-sided\\limits agree};
  \node [barlabel] at ($(constzero)!0.5!(differ)$)
    {one-sided\\limits disagree};

  \begin{pgfonlayer}{background}
    \draw [concept connection, supp!60, ->, >=stealth, shorten >=4pt]
      (cov) edge (known);
  \end{pgfonlayer}
}

\begin{mmpage}
\resizebox{7in}{!}{%
\begin{tikzpicture}[extract levels]
  \pagecanvas
  \techniqueslayer
\end{tikzpicture} 
}
\end{mmpage}

% The same layer, trimmed to its own content.
\begin{mmpage}
\resizebox{2.699in}{!}{%
\begin{tikzpicture}[extract levels]
  \techniqueslayer
\end{tikzpicture} 
}
\end{mmpage}

% ===================================================================
%  Pages 12 and 13 -- the continuity layer: the root and the whole
%  `continuity at a' branch, with the graphs and the IVT note that hang
%  off it on page 1.
%
%  Page 12 is the layer on page 1's canvas, page 13 the same drawing
%  trimmed to its own content.
% ===================================================================
\newcommand{\continuitylayer}{%
  \node [concept, minimum size=3.8cm, text width=3.2cm, font=\small] (root)
    {limit at a constant \(a\)\\[0.8ex]
     % The nested picture inherits the mindmap's node sizing, which flings
     % these two labels centimetres away; reset it here.
     \begin{tikzpicture}[baseline,
         every node/.style={inner sep=1.5pt, minimum size=0pt, minimum width=0pt,
                            minimum height=0pt, text width=, align=center}]
       \draw [->, thin] (0,0) -- (2.5,0) node [right] {\scriptsize\(x\)};
       \fill (1.25,0) circle (1.4pt);
       \node [below] at (1.25,0) {\scriptsize\(a\)};
     \end{tikzpicture}}
    [clockwise from=150]
    child [concept color=gray!20,
           every node/.append style={text width=2.1cm},
           level 2/.append style={level distance=5cm, sibling angle=50},
           level 3/.append style={level distance=3.6cm, sibling angle=52}] { 
      % Wide enough to keep the test on one line -- a node option beats the
      % 2.1cm this subtree sets for everything below it.
      node [minimum size=3.2cm, text width=3.4cm, font=\footnotesize] (conttest)
        {continuity at \(a\)\sub{\lim_{x \to a} f(x) \overset{?}{=} f(a)}}
      [counterclockwise from=75]
      child [concept color=substcase] { 
        node (yes) {continuous}
        [counterclockwise from=65]
        child [concept color=blue!20] { node (ivt) {Intermediate\\Value Theorem} }
        child { node (families) [text width=2.6cm] {can use direct\\substitution\note{to evaluate limits of\\polynomials, rational,\\trigs, \(e^{x}\), \(\ln(x)\)\\in their domain}} }
      }
      child { node (leftcont) [minimum size=2.2cm, text width=2.4cm]
        {left-continuous\note{\(\lim_{x \to a^{-}} f(x) \overset{?}{=} f(a)\)}} }
      child { node (rightcont) [minimum size=2.2cm, text width=2.4cm]
        {right-continuous\note{\(\lim_{x \to a^{+}} f(x) \overset{?}{=} f(a)\)}} }
      child [concept color=gray!20,
             level 3/.append style={level distance=5.2cm, sibling angle=40}] { 
        node (no) {discontinuous\note{calculate both\\one-sided limits}}
        [counterclockwise from=187]
        % Cyan for the two cases where both one-sided limits exist -- not gold,
        % which is the connector-label colour.  The infinite case keeps the
        % pink it shares with nonzero/0.
        child [concept color=bothexist] { node (removable) {removable} }
        child [concept color=bothexist] { node (jump) {jump} }
        child [concept color=pink!50] { node (infinite) [text width=2.6cm]
          {infinite and \(f\) has\\a vertical\\asymptote at \(a\)} }
      }
    }
    ;

  \node [concept color=white, annotation, right, text width=4.2cm,
         font=\scriptsize, align=left] at (ivt.east)
    {IVT applies to continuous function over \([a,b]\) and \(f(a) \ne f(b)\)};

  \node [barlabel] at ($(conttest)!0.5!(yes)$) {yes};
  \node [barlabel] at ($(conttest)!0.5!(no)$) {no};

  \node [barlabel] at ($(no)!0.5!(removable)$)
    {both exist and agree,\\but \(\ne f(a)\)};
  \node [barlabel] at ($(no)!0.5!(jump)$) {both exist\\but disagree};
  \node [barlabel] at ($(no)!0.5!(infinite)$)
    {one of them\\is \(\pm\infty\)};

  % The one-sided continuity pictures, smaller because there are four of them
  % in a crowded corner.  The last is stacked above its sibling rather than
  % anchored straight off the node: everything below rightcont is taken.
  \node [concept color=white, annotation, text width=, above]
    at (leftcont.north) {\includegraphics[width=0.75in, page=13]{plot-limit-intro}};
  \node [concept color=white, annotation, text width=, above left]
    at (leftcont.west) {\includegraphics[width=0.75in, page=14]{plot-limit-intro}};
  \node [concept color=white, annotation, text width=, left]
    at (rightcont.west) {\includegraphics[width=0.75in, page=15]{plot-limit-intro}};
  \node [concept color=white, annotation, text width=, left]
    at ([shift={(0,2.1cm)}]rightcont.west) {\includegraphics[width=0.75in, page=16]{plot-limit-intro}};
}

\begin{mmpage}
\resizebox{7in}{!}{%
\begin{tikzpicture}[map levels]
  \pagecanvas
  \continuitylayer
\end{tikzpicture} 
}
\end{mmpage}

% The same layer, trimmed to its own content.
\begin{mmpage}
\resizebox{4.371in}{!}{%
\begin{tikzpicture}[map levels]
  \continuitylayer
\end{tikzpicture} 
}
\end{mmpage}

% ===================================================================
%  Pages 14 and 15 -- the continuity test: the root, `continuity at a',
%  and the four answers it can have.  What each answer leads to is on the
%  pages around it.
%
%  Page 14 is the layer on page 1's canvas, page 15 the same drawing
%  trimmed to its own content.
% ===================================================================
\newcommand{\continuitytestlayer}{%
  \node [concept, minimum size=3.8cm, text width=3.2cm, font=\small] (root)
    {limit at a constant \(a\)\\[0.8ex]
     % The nested picture inherits the mindmap's node sizing, which flings
     % these two labels centimetres away; reset it here.
     \begin{tikzpicture}[baseline,
         every node/.style={inner sep=1.5pt, minimum size=0pt, minimum width=0pt,
                            minimum height=0pt, text width=, align=center}]
       \draw [->, thin] (0,0) -- (2.5,0) node [right] {\scriptsize\(x\)};
       \fill (1.25,0) circle (1.4pt);
       \node [below] at (1.25,0) {\scriptsize\(a\)};
     \end{tikzpicture}}
    [clockwise from=150]
    child [concept color=gray!20,
           every node/.append style={text width=2.1cm},
           level 2/.append style={level distance=5cm, sibling angle=50},
           level 3/.append style={level distance=3.6cm, sibling angle=52}] { 
      % Wide enough to keep the test on one line -- a node option beats the
      % 2.1cm this subtree sets for everything below it.
      node [minimum size=3.2cm, text width=3.4cm, font=\footnotesize] (conttest)
        {continuity at \(a\)\sub{\lim_{x \to a} f(x) \overset{?}{=} f(a)}}
      [counterclockwise from=75]
      child [concept color=substcase] { node (yes) {continuous} }
      child { node (leftcont) [minimum size=2.2cm, text width=2.4cm]
        {left-continuous\note{\(\lim_{x \to a^{-}} f(x) \overset{?}{=} f(a)\)}} }
      child { node (rightcont) [minimum size=2.2cm, text width=2.4cm]
        {right-continuous\note{\(\lim_{x \to a^{+}} f(x) \overset{?}{=} f(a)\)}} }
      child [concept color=gray!20] { node (no)
        {discontinuous\note{calculate both\\one-sided limits}} }
    }
    ;

  \node [barlabel] at ($(conttest)!0.5!(yes)$) {yes};
  \node [barlabel] at ($(conttest)!0.5!(no)$) {no};

  % The one-sided continuity pictures, smaller because there are four of them
  % in a crowded corner.  The last is stacked above its sibling rather than
  % anchored straight off the node: everything below rightcont is taken.
  \node [concept color=white, annotation, text width=, above]
    at (leftcont.north) {\includegraphics[width=0.75in, page=13]{plot-limit-intro}};
  \node [concept color=white, annotation, text width=, above left]
    at (leftcont.west) {\includegraphics[width=0.75in, page=14]{plot-limit-intro}};
  \node [concept color=white, annotation, text width=, left]
    at (rightcont.west) {\includegraphics[width=0.75in, page=15]{plot-limit-intro}};
  \node [concept color=white, annotation, text width=, left]
    at ([shift={(0,2.1cm)}]rightcont.west) {\includegraphics[width=0.75in, page=16]{plot-limit-intro}};
}

\begin{mmpage}
\resizebox{7in}{!}{%
\begin{tikzpicture}[map levels]
  \pagecanvas
  \continuitytestlayer
\end{tikzpicture} 
}
\end{mmpage}

% The same layer, trimmed to its own content.
\begin{mmpage}
\resizebox{3.712in}{!}{%
\begin{tikzpicture}[map levels]
  \continuitytestlayer
\end{tikzpicture} 
}
\end{mmpage}

% ===================================================================
%  Pages 16 and 17 -- why substitution works: `continuity at a' and what
%  a continuous function buys you, beside the one case of direct
%  substitution that needs nothing further, and the arrow from that case
%  to continuity.
%
%  Two subtrees with different level distances share this picture, so each
%  is drawn in a scope of its own rather than the levels being set on the
%  picture.
%
%  Page 16 is the layer on page 1's canvas, page 17 the same drawing
%  trimmed to its own content.
% ===================================================================
\newcommand{\substitutionworkslayer}{%
  % `missing' children hold the places of the siblings this layer leaves
  % out, so the ones it keeps fan out at page 1's angles without the start
  % of the fan having to be recomputed.
  \begin{scope}[continuity levels]
  % `conttest' is this layer's root, placed where page 1 puts it: the first
  % of the root's branches, 150 degrees out at the level-1 5.5cm.  The
  % styling is what `every child node' gives it there, spelled out because
  % it is no longer a child of anything.
  \node [concept color=gray!20, concept, minimum size=3.2cm, text width=3.4cm,
         align=center, font=\footnotesize] (conttest) at (150:5.5cm)
    {continuity at \(a\)\sub{\lim_{x \to a} f(x) \overset{?}{=} f(a)}}
    [counterclockwise from=75]
    child [concept color=substcase] { 
      node (yes) {continuous}
      [counterclockwise from=65]
      child [concept color=blue!20] { node (ivt) {Intermediate\\Value Theorem} }
      child { node (families) [text width=2.6cm] {can use direct\\substitution\note{to evaluate limits of\\polynomials, rational,\\trigs, \(e^{x}\), \(\ln(x)\)\\in their domain}} }
    }
    ;
  \end{scope}

  \begin{scope}[technique levels]
    \node [concept color=gray!20, concept, minimum size=2.4cm, text width=2.2cm,
           align=center, font=\footnotesize] (dirsub) at (30:5.5cm)
      {try direct\\substitution\note{is the denominator \(0\)?}}
      [counterclockwise from=-30]
      child [missing]
      child [missing]
      child [concept color=substcase] { node (algno) [text width=2.1cm]
          {\(\frac{\text{any number}}{\text{nonzero}}\)} }
      ;
  \end{scope}

  \node [barlabel] at ($(conttest)!0.5!(yes)$) {yes};
  \node [concept color=white, annotation, right, text width=4.2cm,
         font=\scriptsize, align=left] at (ivt.east)
    {IVT applies to continuous function over \([a,b]\) and \(f(a) \ne f(b)\)};

  % The two decision trees reach the same place by different routes.
  \begin{pgfonlayer}{background}
    \draw [concept connection, substline, ->, >=stealth, shorten >=4pt]
      (algno) edge (yes);
  \end{pgfonlayer}
}

\begin{mmpage}
\resizebox{7in}{!}{%
\begin{tikzpicture}[mindmap base]
  \pagecanvas
  \substitutionworkslayer
\end{tikzpicture} 
}
\end{mmpage}

% The same layer, trimmed to its own content.
\begin{mmpage}
\resizebox{3.092in}{!}{%
\begin{tikzpicture}[mindmap base]
  \substitutionworkslayer
\end{tikzpicture} 
}
\end{mmpage}

% ===================================================================
%  Pages 18 and 19 -- the discontinuous branch: `continuity at a', the
%  answer that ends the question, and the three ways a function can fail
%  to be continuous.
%
%  Page 18 is the layer on page 1's canvas, page 19 the same drawing
%  trimmed to its own content.
% ===================================================================
\newcommand{\discontinuouslayer}{%
  \begin{scope}[continuity levels]
  % `conttest' is this layer's root, placed where page 1 puts it: the first
  % of the root's branches, 150 degrees out at the level-1 5.5cm.  The
  % styling is what `every child node' gives it there, spelled out because
  % it is no longer a child of anything.
  \node [concept color=gray!20, concept, minimum size=3.2cm, text width=3.4cm,
         align=center, font=\footnotesize] (conttest) at (150:5.5cm)
    {continuity at \(a\)\sub{\lim_{x \to a} f(x) \overset{?}{=} f(a)}}
    [counterclockwise from=75]
    child [concept color=substcase] { node (yes) {continuous} }
    % The two one-sided tests sit between them on page 1; `missing' keeps
    % their places so `no' still lands on its own bar.
    child [missing]
    child [missing]
    child [concept color=gray!20,
           level 2/.append style={level distance=5.2cm, sibling angle=40}] { 
      node (no) {discontinuous\note{calculate both\\one-sided limits}}
      [counterclockwise from=187]
      % Cyan for the two cases where both one-sided limits exist -- not gold,
      % which is the connector-label colour.  The infinite case keeps the
      % pink it shares with nonzero/0.
      child [concept color=bothexist] { node (removable) {removable} }
      child [concept color=bothexist] { node (jump) {jump} }
      child [concept color=pink!50] { node (infinite) [text width=2.6cm]
        {infinite and \(f\) has\\a vertical\\asymptote at \(a\)} }
    }
    ;
  \end{scope}

  \node [barlabel] at ($(conttest)!0.5!(yes)$) {yes};
  \node [barlabel] at ($(conttest)!0.5!(no)$) {no};
  \node [barlabel] at ($(no)!0.5!(removable)$)
    {both exist and agree,\\but \(\ne f(a)\)};
  \node [barlabel] at ($(no)!0.5!(jump)$) {both exist\\but disagree};
  \node [barlabel] at ($(no)!0.5!(infinite)$)
    {one of them\\is \(\pm\infty\)};
}

\begin{mmpage}
\resizebox{7in}{!}{%
\begin{tikzpicture}[mindmap base]
  \pagecanvas
  \discontinuouslayer
\end{tikzpicture} 
}
\end{mmpage}

% The same layer, trimmed to its own content.
\begin{mmpage}
\resizebox{2.949in}{!}{%
\begin{tikzpicture}[mindmap base]
  \discontinuouslayer
\end{tikzpicture} 
}
\end{mmpage}

\end{document}
